\documentclass{article}
\usepackage[margin=0.5in]{geometry}
\usepackage{amsmath, amssymb, booktabs}
\pagestyle{empty}
\begin{document}

In the Ordinary Least Squares Problem, we either assumed no noise or treated noise as additive factor on the underlying model output function.
In a Generalized Linear Model, we treat the model output function as the mean (expected value) of a distribution (for each observation at a time).

\section{How The Gauss Normal Equations Are The Result Of A Normal Distributed Maximum Likelihood Estimator}
Assume each observation is generated by the model $y_i = a_i^\top b + \varepsilon_i$ with target $y_i$, input $a_i$, parameters $b$, and
independent, identically distributed Gaussian noise $\varepsilon_i \sim \mathcal{N}(0, \sigma^2)$. The likelihood of observing $y_i$ given $a_i$ and $b$ follows the Gaussian probability density function
\[
    p(y_i \mid a_i, b) = \frac{1}{\sqrt{2\pi\sigma^2}}
    \exp\!\left(-\frac{(y_i - a_i^\top b)^2}{2\sigma^2}\right).
\]
The joint likelihood $\mathcal{L}$ over all $m$ independent observations is a function of our assumed parameters $b$:
\[
    \mathcal{L}(b) = \prod_{i=1}^{m} p(y_i \mid a_i, b)
    = \prod_{i=1}^{m} \frac{1}{\sqrt{2\pi\sigma^2}}
      \exp\!\left(-\frac{(y_i - a_i^\top b)^2}{2\sigma^2}\right).
\]
Taking the log turns the product into a sum $(\log\left(\frac{1}{\sqrt{2\pi\sigma^2}}\right) = \log\left((2\pi\sigma^2)^{-1/2}\right) = -\frac{1}{2}\log(2\pi\sigma^2)
)$:
\[
    \log \mathcal{L}(b)
    = \sum_{i=1}^{m} \log \frac{1}{\sqrt{2\pi\sigma^2}}
      + \sum_{i=1}^{m} \log \exp\!\left(-\frac{(y_i - a_i^\top b)^2}{2\sigma^2}\right)
    = -\frac{m}{2}\log(2\pi\sigma^2)
      - \frac{1}{2\sigma^2}\sum_{i=1}^{m}(y_i - a_i^\top b)^2
    = \mathrm{c} - \frac{1}{2\sigma^2}\|y - Ab\|^2.
\]
Only the $b$-dependent term affects the maximum; its negative sign makes maximising the log-likelihood equivalent to minimising the squared residual:
\[
    \underset{b}{\arg\max}\ \log \mathcal{L}(b)
    \;=\;
    \underset{b}{\arg\min}\ \|y - Ab\|^2.
\]
Hence the least-squares loss is equivalent to the Gaussian log-likelihood, and the Normal Equation $A^\top A b = A^\top y$
is its maximum-likelihood condition.

\section{Other Distributions: Towards Generalized Linear Models}
The Gaussian derivation generalises. Replace the Gaussian assumption on $y_i$ with any other distribution and again maximise $\log \mathcal{L}(b)$. The two things that change are: (1) the distribution of $y_i$, and (2) how the linear predictor $\eta_i = a_i^\top b$ connects to the mean $\mu_i = \mathbb{E}[y_i]$. In the Gaussian case this connection was simply the identity $\mu_i = \eta_i$. In general, that connection is the \emph{link function} $g$, defined by $g(\mu_i) = \eta_i$, so that $\mu_i = g^{-1}(\eta_i)$.

\vspace{3mm}
A link function $g$ is called \emph{canonical} if it arises naturally from writing the distribution of $y_i$ in \emph{exponential family} form:
\[
    f(y_i \mid \theta_i) = h(y_i)\exp\Big(\theta_i y_i - A(\theta_i)\Big), \quad \text{and to simplify later derivations we use} \quad \log f =  \log h(y_i) + \theta_i y_i - A(\theta_i).
\]
where $\theta_i$ is the \emph{natural parameter}, $A(\theta_i)$ is the \emph{log-partition function} (normalising the distribution so it integrates to 1), and $h(y_i)$ is a base measure depending only on $y_i$.
The canonical link is the choice that sets $\theta_i = \eta_i$, under which the gradient of the log-likelihood takes the same form as in ordinary linear regression:
\[
    \nabla_b \log \mathcal{L}(b) = A^\top(y - \mu).
\]

\subsection{Bernoulli Distribution}
For Bernoulli targets $y_i \in \{0,1\}$, the mean is the probability of success, $\mu_i \in (0,1)$, with $p(y_i \mid \mu_i) = \mu_i^{y_i}(1-\mu_i)^{1-y_i}$ as density. Taking the log and using $\log a - \log b = \log(a/b)$:
\[
\log p(y_i \mid \mu_i)
= y_i \log \mu_i + (1-y_i)\log(1-\mu_i) = y_i \log\frac{\mu_i}{1-\mu_i} + \log(1-\mu_i).
\]
The logarithmized exponential family form is $\log(h(y_i)) + \theta_i y_i - A(\theta_i)$. Clearly, $\log(h(y_i)) = 0 \implies h(y_i) = 1$. The coefficient of $y_i$ in $\log p$ is $\log\frac{\mu_i}{1-\mu_i}$, so we read off $\theta_i = \log\frac{\mu_i}{1-\mu_i} \Leftrightarrow \mu_i = \frac{1}{1+e^{-\theta_i}}$. 
The remaining term must therefore equal $-A(\theta_i)$. Substituting as functions of $\theta_i$:
\[
\log p(y_i \mid \theta_i) = y_i\theta_i - \log(1+e^{\theta_i}),
\]
The canonical link sets $\theta_i = \eta_i = a_i^\top b$, giving
\[
\log\frac{\mu_i}{1-\mu_i} = a_i^\top b \;\implies\; \mu_i = \frac{1}{1+e^{-a_i^\top b}}.
\]

Substituting $\theta_i = a_i^\top b$ into $\log p(y_i \mid \theta_i)$, each observation contributes
$y_i(a_i^\top b) - \log(1 + e^{a_i^\top b})$ to the log-likelihood. Summing over all $m$ independent
observations (as we already logarithmized everything):
\[
    \log \mathcal{L}(b) = \sum_{i=1}^{m} \left[y_i(a_i^\top b) - \log(1 + e^{a_i^\top b})\right].
\]
Differentiating term by term: $\nabla_b(a_i^\top b) = a_i$, and by the chain rule
$\nabla_b \log(1 + e^{a_i^\top b}) = \frac{e^{a_i^\top b}}{1+e^{a_i^\top b}}\,a_i = \mu_i a_i$.
Therefore:
\[
    \nabla_b \log \mathcal{L}(b) = \sum_{i=1}^{m} a_i(y_i - \mu_i) = A^\top(y - \mu),
\]
where $A$ is the design matrix with rows $a_i^\top$ and $\mu = (\mu_1, \dots, \mu_m)^\top$.

Although the Bernoulli log-likelihood is not quadratic, its gradient has the same form $A^\top(y-\mu)$
as in the Gaussian case.

\subsection{Gamma Distribution}
For positive continuous targets $y_i > 0$, the Gamma distribution with known shape $k > 0$
and mean $\mu_i > 0$ has density
\[
    p(y_i \mid \mu_i, k) = \frac{1}{\Gamma(k)}\!\left(\frac{k}{\mu_i}\right)^{\!k}
    y_i^{k-1} \exp\!\left(-\frac{k\, y_i}{\mu_i}\right).
\]
Taking the log:
\[
    \log p(y_i \mid \mu_i, k)
    = \underbrace{\bigl[(k-1)\log y_i - \log\Gamma(k) + k\log k\bigr]}_{\log h(y_i)}
    + \underbrace{\left(-\frac{k}{\mu_i}\right)}_{\theta_i} y_i
    - \underbrace{\bigl(-k\log(-\theta_i)\bigr)}_{A(\theta_i)}.
\]
Matching against the exponential family log-form $\log h(y_i) + \theta_i y_i - A(\theta_i)$,
we read off $\theta_i = -k/\mu_i$, giving canonical link $g(\mu_i) = -k/\mu_i$ (the inverse link).
In practice the \emph{log link} $g(\mu_i) = \log\mu_i = a_i^\top b$ is used instead,
since it guarantees $\mu_i = e^{a_i^\top b} > 0$ for any $b$ and gives
coefficients a multiplicative (percentage) interpretation.

Substituting $\mu_i = e^{a_i^\top b}$, dropping terms constant in $b$ and applying the chain rule $\nabla_b\, y_i e^{-a_i^\top b} = -y_i/\mu_i \cdot a_i$:
\[
    \log \mathcal{L}(b)
    = \mathrm{const} - k\sum_{i=1}^{m}\Bigl[a_i^\top b + y_i\,e^{-a_i^\top b}\Bigr], \quad \nabla_b \log \mathcal{L}(b)
    = k\,A^\top\!\left(\frac{y}{\mu} - \mathbf{1}\right), \quad \mu_i = e^{a_i^\top b}.
\]
    
where $y/\mu$ is elementwise division. The residual is now relative: $y_i/\mu_i - 1$
rather than $y_i - \mu_i$, reflecting that Gamma penalises proportional errors
($\mathrm{Var}(y_i) = \mu_i^2/k$ grows with the mean), which is appropriate for
ecological measurements spanning orders of magnitude.

Note: The canonical link is a natural choice arising from the exponential family form, but any correct monotone differentiable mapping is valid link. 
For the Gamma distribution the canonical link is avoided in practice since it can change the sign for arbitrary b, but the Gamma distribution's mean is always positive.

\section{Comparison}
\begin{table}[h]
\centering
\begin{tabular}{@{}llllll@{}}
\toprule
\textbf{Noise} & \textbf{Distribution on } $y_i$ & \textbf{Link} $g(\mu_i)$ & \textbf{Mean} $\mu_i$ & \textbf{Loss / Log-likelihood} & \textbf{Solution} \\
\midrule
None
  & $y_i = a_i^\top b$ (exact)
  & Identity
  & $a_i^\top b$
  & $\|y - Ab\|^2 = 0$
  & $A^\top A b = A^\top y$, direct (QR) \\[4pt]
Gaussian
  & $y_i \sim \mathcal{N}(a_i^\top b,\;\sigma^2)$
  & Identity
  & $a_i^\top b$
  & $-\tfrac{1}{2\sigma^2}\|y-Ab\|^2$\;
  & $A^\top A b = A^\top y$, direct (QR) \\[4pt]
Gamma
  & $y_i \sim \mathrm{Gamma}(k,\;\mu_i/k)$
  & Log
  & $e^{a_i^\top b}$
  & $-k\sum_i\bigl[a_i^\top b + y_i e^{-a_i^\top b}\bigr]$\;
  & $\nabla_b \log\mathcal{L}=0$ (grad descent) \\
\bottomrule
\end{tabular}
\caption{Comparison of the three regression settings. The no-noise and Gaussian cases share the same normal equation because the Gaussian log-likelihood is quadratic in $b$, making its gradient linear and exactly solvable. The Gamma log-likelihood contains $e^{-a_i^\top b}$, breaking this structure and requiring iterative optimisation.}
\end{table}

\end{document}